\documentclass[12pt,a4paper]{article}
\usepackage[UTF8]{ctex}        % 中文支持
\usepackage{geometry}
\geometry{left=2.54cm,right=2.54cm,top=3.18cm,bottom=3.18cm}
\usepackage{enumerate}        % 列表格式
\usepackage{titlesec}         % 标题格式
\usepackage{xcolor}           % 颜色包
\usepackage{indentfirst}
\setlength{\parindent}{2em}
\usepackage{enumitem}
\usepackage{enumerate}

% 字体设置


% ===================== 自定义功能区 =====================
\newcommand{\textred}[1]{\textcolor{red}{#1}}           %红色字体
\newcommand{\textgreen}[1]{\textcolor{green}{#1}}       %绿色字体
\newcommand{\textblue}[1]{\textcolor{blue}{#1}}         %绿色字体
\newcommand{\textyellow}[1]{\textcolor{yellow}{#1}}     %黄色字体
% \newcommand{\wordnote}[2]{#1\,\raisebox{0.25ex}{\scriptsize #2}}
\newcommand{\upmean}[1]{\raisebox{0.25ex}{\scriptsize #1}} %标注





%========================== 正文 ============================


\begin{document}


% 书籍《我心态超好》
\section*{《我心态超好》}




\subsection*{第一部分 第二章 “夜不能寐”是身体发出的危险信号 }

以下几类人很容易过度紧张找上门，应时刻注意。他们分别是：一丝不苟的人、较为敏感的人
和追求完美的人。

“闲余时间”的消失并非是引起过度紧张的唯一原因。

SNS(social network serve)的广泛使用也是引起过度焦虑和紧张的一个因素。

如果这些智能设备\upmean{电脑、手机}已经成为你生活中不可或缺的一部分，
并且占据了你大部分的时间，那么你要注意了，它们可能正在悄悄吞噬你的身心健康。

VDT综合征（Visual Display Terminal）：是指因长时间面对电脑、手机等
智能设备，引起眼疲劳、肩痛、头痛、腰痛、倦怠感、眩晕等不适症状。

对电脑等智能设备的使用时间和方法给出了指导性标准。
\begin{enumerate}
    \item 每次工作时间不超过1h，每隔1h休息10~15min。
    \item 根据实际的工作请款安排小憩时间，避免连续工作。
    \item 为了预防腰背酸痛，请勿长时间保持一个姿势。
    
    另外，如果在午休和上下班的路上尽可能不看手机，也能预防VDT综合征，
    缓解过度紧张。
\end{enumerate}

若想应对过度紧张和VDT综合征带来的各项危害，最简单的方法是：让副交感神经占据上风，
有意识地创造放松空间。

我将让副交感神经兴奋起来的时间称为“放松时间”。对于那些容易紧张焦虑的人，
可以按照下面的方法制造“放松时间”。
\begin{enumerate}
    \item 制定只属于自己的“睡眠日”
    
    因此一旦发觉过度紧张焦虑来到了你的身边，首先要在制定属于自己的“睡眠日”。
    “今天我们的任务就是好好睡觉！”一旦下定决心，就要尽力实现。不着急的工作
    不妨暂时搁置，家务活也不必强求。总之，天大地大睡觉最大。

    \item 回家后，远离手机
    
    SNS或者手机游戏会让焦虑情绪雪上加霜。
    尽可能不看手机是制造“放松时间”的前提。

    \item 想要和过度紧张抗衡，就要学会放下一切，悠哉度日。
    
    总之，伸展筋骨、放松心情是缓解焦虑的精髓所在。

    \item 终止那些容易带来焦虑的学习
    
    学习这件事，和工作一样，也能带来紧张、焦虑感。因此，
    一旦出现过度紧张的症状，及时终止方为上策。

    总而言之，神经过度紧张时，首先要考虑用睡觉来缓解，
    因此回到家后，切记不要晚睡。

    \item 休息日不约见，不远行
    
    在周末，如果感到身心俱疲、精神紧张，切记不要和让
    你感到拘谨的人见面，也最好原理人群喧嚣之地，并且
    尽可能地不要远行。

    只需按照自己地节奏、悠闲自得地度过就好。

    如果非常想外出游玩，可以选择近一点地地方，相关日程
    也不必安排得非常紧凑。随心所欲地游玩何尝不是一件
    美事呢！

\end{enumerate}





\subsection*{第三部分 第十章 睡眠 }

压力会首先侵蚀我们地睡眠。

正在阅读本书的你是否也还清了所有的睡眠负债呢？如果出现
下述症状，证明你还处于“欠债”状态。

\begin{enumerate}
    \item 三餐后或者坐地铁上总是犯困，经常睡着。
    \item 白天如果不抽烟、不喝咖啡，就会变得萎靡不振。
    \item 在枯燥的会议上经常犯困，处于似睡非睡状态。
    \item 晚上一上床倒头就睡，可谓“一秒入睡”。
    \item 在开车的过程中，只要车一听（等待红绿灯期间等），困意就席卷而来。
\end{enumerate}\vspace{-15pt}
{\raggedleft \scriptsize  超过中午12点的补交，会严重扰乱人体的生物钟。\par}

无论多么热衷于工作，一定要劳逸结合，否则就会积劳成疾。

{\centering {\Large {\kaishu \textcolor{gray}{调理睡眠十分重要，六大法宝轻松解决}}}\par}


为提高失眠质量，睡觉前的放松必不可少。

人不是机器，想要在睡前保持内心的平静是需要一定时间的。

以下六大秘诀能轻松帮你解决睡眠问题。

\begin{enumerate}
    \item 睡前1-2 h不碰手机、电脑等电子产品
    
    电子产品发出的蓝光会让大脑误以为现在是白天，当然不会产生倦意。此外，卧室灯不要使用
    白炽灯，因为白炽灯发出的白光波长和蓝光相似。

    \item 通过享受晚餐、和家人的团聚、泡澡来放松身心
    
    但不建议睡觉前用温度较高的水长时间泡澡。

    \item 晚上不要摄入含咖啡因的饮品
    
    \item 晚饭和睡觉只要间隔2~3 h
    
    \item 睡前3 h不饮酒
    
    \item 在安静黑暗、温度适宜的卧室睡觉
    
    声音和光亮不利于熟睡。这样的睡眠状态让人积存的疲劳感很难得到缓解。

    \footnotesize \textcolor{teal}{注:好的睡眠可以增强人体免疫力，
    远离感冒等疾病的入侵，也可以提高工作效率，减少犯错。}

\end{enumerate}





\subsection*{第三部分 第十一章 饮食 }

{\centering {\Large {\kaishu \textcolor{gray}{每顿饭都草草了事的话，和以缓解疲劳？}}}\par}

拒绝敷衍了事的饮食

长期的饮食不规律、不健康，会导致营养元素减少，从而造成营养不良。

众所周知，血清素和多巴胺等神经递质对人体来说至关重要。如果人体营养不良，就意味着没有养料
再提供给这些神经递质了。紧接着，人体就会出现各种症状。

生长激素的分泌和睡眠息息相关，人只有在睡着了，大脑垂体才会分泌生长激素。

健康的诀窍：红：黄：绿 = 1：1：1

\textred{红色食材}：肌肉、血液、内脏所需要的蛋白质。

包括：肉类、鱼类、鸡蛋；豆腐、纳豆；芝士、纯牛奶、酸奶等。

\textyellow{黄色食材}：碳水化合物、食用油等人体能量的来源。

包括：米饭、面食、糕点、薯类；黄油、色拉油、橄榄油、鲜奶油等。

\textgreen{绿色食材}：调节身体各项机能的维生素、矿物质。

包括：蔬菜、藻类、水果等。

工作中每半个小时就起身活动一次。

步行、跑步、骑行都属于有氧运动，如果能长期坚持，
不仅能够强身健体，还能让我们远离精神类疾病。
步行、慢跑、骑行等有一定节奏感的运动会促进大脑分泌血清素。

血清素可以有效对抗压力，而且能让人镇静，减少急躁情绪，带来愉悦感和幸福感。

有田秀穗教授曾说：“一旦开始一定节奏的运动，5分钟后，大脑中的血清素浓度就会有明显上升，20~30分钟后达到峰值。”

阳光也能促进血清素的分泌，因此最好在户外进行慢跑、快走等运动。



\subsection*{第三部分 第十三章 心理 }

{\centering {\Large {\kaishu \textcolor{gray}{内心能量不足时要及时补充}}}\par}

养护心灵最佳方法：给自己直面内心的时间。

条件允许的话，希望大家每天都能挤出点时间直面最真实的自己，倾听自己内心最深处的声音。

如何给内心充电？

\begin{enumerate}
    \item 做一些内心真正期望的，让自己兴奋不已的事情。
    \item 内心平静、身心放松、休息。
    \item 做一些有益健康、心情愉悦的事。
    \item 通过参加一些和自己价值观相符的活动、内心变得充实无比。
\end{enumerate}

身体和心理是密不可分的一个整体，当心理较为脆弱时，身体也会表现出疲惫和不适。当我们察觉到压力发出的信号时，切不可麻痹大意，一定要在第一时间调整睡眠和饮食，让疲劳无机可乘。

摄入营养丰富的食物、保证充足的睡眠，能让内心的能量得到快速恢复。

当发现内心的能量出现下滑时，尽量不要打破已有的生活方式，让生活出现变化。当压力发出信号时，应尽量避免变化的发生。例如：换工作、学习新知识、挑战减肥和戒烟、去远方旅行等，或多或少都会给人带来压力。

当压力发出信号时，应尽量遵从自己的本心，优先考虑“我想···”，想方设法减少“我必须···”的情况发生。

与压力巧妙共存的3R法则：

Reest: 休息和睡眠

Relaxation：放松

Recreation：娱乐




\section*{《认知天性》}

美国前国防部部长唐纳德·拉姆斯菲尔德：

有些事情是已知的已知；有些事情是已知的未知。但是还有未知的未知。








\section*{git基础操作}
git初始化（仅第一次需要）：\par
git init\par
将所有文件放到等待区：\par
全选文件：\par
git add .\par
单个文件：\par
git add name\par
本地打包提交：\par
git commit -m "本次修改说明"\par
绑定远程仓库：\par
git remote add origin https://github.com/username/仓库名.git\par
推送到GitHub云端：\par
git push -u origin main

查看当前绑定的远程仓库：\par
git remote -v\par

情况1：换一个新 GitHub 仓库推送\par
删除旧的远程绑定：\par
git remote remove origin\par
绑定新仓库链接：\par
git remote add origin 新仓库的github链接\par

情况 2：同时存多个仓库（一键推送到两个仓库）：\par

新增第二个远程，取名other\par
git remote add other https://github.com/用户名/第二个仓库.git\par

推送到第一个仓库\par
git push origin main\par
推送到第二个仓库\par
git push other main\par

查看所有提交记录：\par
git log\par


创建分支：\par
git branch \par
创建新分支（两种写法任选其一）:\par
git branch dev\par
创建分支并立刻切换到新分支（最常用）:\par
git checkout -b dev\par
Git 2.23+新版本推荐写法：\par
git switch -c dev\par

切换已有分支：\par
git checkout main\par
删除本地无用分支：\par
git branch -d dev\par
删除github远程分支：\par
git push origin --delete dev\par
拉取远程新建分支到本地\par
git fetch origin\par
git checkout dev\par


\subsection*{以下书籍是为了更加深入地学习git的详细内容}
Git 权威学习资源大全（分 4 类：官方指南书、在线文档、交互式实操、本地自带手册）

1. 在线阅读地址（官方正版简体中文）

https://git-scm.com/book/zh/v2

二、Git 官网官方命令参考手册（查指令专用）

地址：https://git-scm.com/doc

包含所有 Git 命令完整语法，每一条指令的全部参数、细节用法；

搜索框直接输入 branch / push / remote，就能查到这条命令所有隐藏用法；

配套官方简短入门教程：https://git-scm.com/docs/gittutorial

三、GitHub 官方配套文档（专门讲 GitHub 平台操作）

地址：https://docs.github.com/zh

如果你想搞懂：新建仓库、SSH 密钥、PR 合并、代码评审、团队分支规范、GitHub Pages；

解决本地 Git 和 GitHub 网页联动的所有问题。

四、本地终端自带手册（不用联网，随时查）

安装 Git 后，直接在终端输入命令看文档：

\# 查看Git总帮助目录

git help -g

\# 打开官方入门教程

git help tutorial

\# 单独查某一条指令完整文档，比如branch、push

git help branch

git help push

git help remote

五、交互式可视化实操教程（练分支最好玩，直观理解）

Learn Git Branching（分支可视化神器）

地址：https://learngitbranching.js.org/

网页游戏式教学，图形画出分支指针、合并、rebase 全过程；

专门攻克新手最难懂的分支、合并、变基，边敲指令边看动画。

六、进阶实战补充资料

Git Flight Rules：踩坑急救手册（所有撤销代码、冲突修复、丢失版本找回方案）

https://github.com/k88hudson/git-flight-rules

Atlassian Git 教程（团队协作实战）

https://www.atlassian.com/git/tutorials

速查表：GitHub 官方 Git 命令速记图，打印随身看

https://education.github.com/git-cheat-sheet
















\newpage

Shoot for the moon. Even if you miss, you'll land among the stars.

瞄准月亮吧。即使你没射中，也会落在星辰之间。










\end{document}